% FILE: 3-poster/parts/macros.tex
% ==================================================================================================
% Reusable toolkit: type sizes, math macros, layout grid, coordinate helpers,
% card / heading / icon / citation commands, and TikZ arrow styles.
% Requires colors.tex and fonts.tex to be loaded first.

% ===== Type-size hierarchy (absolute points; poster is at true CCN scale) =========================
% Readability floor = the Approach body size (30 pt): no text, INCLUDING inside figures, is set
% below \labelsize. Leading is generous to aerate the page.
% Leading = 140% of the font size (UX default), applied uniformly.
\newcommand{\titlesize}{\fontsize{104}{146}\selectfont}
\newcommand{\authorsize}{\fontsize{54}{76}\selectfont}
\newcommand{\affilsize}{\fontsize{42}{59}\selectfont}
\newcommand{\contactsize}{\fontsize{42}{59}\selectfont}
\newcommand{\headinglabelsize}{\fontsize{62}{87}\selectfont}
\newcommand{\headingdescsize}{\fontsize{48}{67}\selectfont}
\newcommand{\resultsize}{\fontsize{64}{90}\selectfont}
\newcommand{\takeawaysize}{\fontsize{54}{76}\selectfont}
\newcommand{\bodysize}{\fontsize{40}{56}\selectfont}
\newcommand{\bodybigsize}{\fontsize{52}{70}\selectfont}   % prominent formulas
\newcommand{\smallsize}{\fontsize{40}{56}\selectfont}     % floor (40 pt)
\newcommand{\labelsize}{\fontsize{40}{56}\selectfont}     % floor, used inside figures
\newcommand{\fooLabelsize}{\fontsize{42}{59}\selectfont}
\newcommand{\minorlabelsize}{\fontsize{24}{28}\selectfont} % minor in-figure labels (axis ticks)

% ===== Math macros ================================================================================
% Vectors use \vec{} (arrow); bold is reserved for matrices (J, I, Sigma).
\newcommand{\vk}{\vec{\kappa}}
\newcommand{\vu}{\vec{u}}
\newcommand{\vxi}{\vec{\xi}}
\newcommand{\veta}{\vec{\eta}}
\newcommand{\vv}{\vec{v}}
\newcommand{\mI}{\mathbf{I}}
\newcommand{\mJ}{\mathbf{J}}
\newcommand{\vm}{\vec{m}}
\newcommand{\vn}{\vec{n}}
\newcommand{\Sig}{\boldsymbol{\Sigma}}
% In-math text label rendered in the math (serif) text font, so words match the symbols.
\newcommand{\mlab}[1]{\text{{\mathtextfont #1}}}
% An \underbracket's READING, pinned to 32 pt INDEPENDENTLY of the formula's size. \mlab would follow
% the surrounding math size, so enlarging a formula for emphasis would enlarge its gloss with it -- and
% the gloss is not part of the statement: it is the words that read the statement, and it must stay
% quiet. Use \brlab (not \mlab) under a brace whenever the formula is set above body size.
\newcommand{\brlab}[1]{\text{{\mathtextfont\fontsize{32}{36}\selectfont #1}}}
% Inline math sized to visually match 40 pt sans body text, for use INSIDE TikZ nodes (where neither
% \DeclareMathSizes nor a local \fontsize reaches the math size). \scalebox enlarges the rendered
% symbols physically about the baseline, so serif states/matrices read at the surrounding prose size.
\newcommand{\tmath}[1]{\scalebox{1.4}{\ensuremath{#1}}}
\newcommand{\RR}{\mathbb{R}}
% multi-letter "mix" index set one math style smaller, to read like the 1/2 indices
\newcommand{\kmix}{\vk_{\scriptscriptstyle\mathrm{mix}}}
\newcommand{\umix}{\vu_{\scriptscriptstyle\mathrm{mix}}}
\newcommand{\kref}{\vk_{\scriptscriptstyle\mathrm{ref}}}
\newcommand{\Reff}{R_{\mathrm{eff}}}
\newcommand{\Cspan}{\mathcal{C}_{12}}

% ===== Layout grid (cm; origin at the page's south-west corner) ===================================
% Page (true CCN 2026 portrait size).
\def\PW{147.32}\def\PH{172.72}
% Outer margins.
\def\Mside{5}\def\Mtop{4}\def\Mbot{3}
% Columns (x). The two content columns stack their blocks INDEPENDENTLY (no rigid 2x2),
% which lets each side flow to fit its content.
\def\gXl{5}            % left content edge
\def\gXr{142.32}       % right content edge (PW - Mside)
\def\gWl{65}           % left column width (shared by ALL left-column cards; NOT to be changed for lens-figure sizing -- size the lens plots via \lensscale in block-starting.tex instead)
\def\gWr{68.32}        % right column width
\def\gXrcL{74}         % right column, left edge
\def\gWfull{137.32}    % full content width (footer / supporting)

% Region map (cm; y = NORTH edge of each block, h = height). Single source of truth.
%   FULL width : header banner (spans \gWfull, like the footer) -- BOTH columns start below it.
\def\yHead{168.72}                           % 168.72 .. 159.63 (banner is 9.09 cm tall; the build
                                             % logs its measured height as "[layout] header banner
                                             % height" -- re-read it after any header edit)
% BOTH columns hang from the SAME top edge, 157 -- a ~2.6 cm gap below the banner (the columns'
% internal card gaps are ~1 cm, so this reads as a section break, not as dead space). Changing that
% top edge means re-deriving BOTH chains below: each is rigid (a card's y is its predecessor's bottom
% minus the gap), and each is pinned at the far end (the left column must still bottom out at 8.5,
% just above the footer divider at y=7.3).
%   LEFT column : starting (now INCLUDING the open question) -> approach
% Both card heights are DERIVED, not guessed: each block boxes its content and logs it as
% "[layout] <block> content height"; the card needs that + 22 mm (the cardbase top+bottom padding).
% Re-read those log lines after editing either card -- a tcolorbox with valign=top spills BELOW its
% frame with NO overfull-vbox warning, so an undersized height is invisible in the log.
% THE LEFT COLUMN IS FULL, to within ~3 mm. Restoring the cards' 56 pt body leading (see the
% `fontupper' note in block-starting.tex) cost ~10 cm, paid for by tightening the section gaps, the
% modality images, the badges, and the archetype plots (\archscale). ANY growth in either card must
% now be funded by an explicit cut in the other.
\def\yStart{158.6}\def\hStart{62.0}          % 158.6 .. 96.6 (content 59.7 + 2.2 padding, 0.1 spare).
                                             % The Question is EMBEDDED at the bottom of this card;
                                             % \yQuest / \hQuest are retired.
                                             % GREW 61.5 -> 62.0 for the FOURTH archetype: the row
                                             % gained a citation line under each badge (+2.6 cm) and a
                                             % two-line x-label (+1.2 cm), while narrowing the plots
                                             % four-across returned 2.4 cm. Net +0.4, funded below.
\def\yAppr{95.6}\def\hAppr{87.6}             % 95.6 .. 8.0 (1 cm gap above; ~0.7 cm above the footer divider at y=7.3). Content measures 82.9 -> ~4.7 cm of slack (was 5.2: the Starting card's fourth archetype spent 0.5 of it, per the "any growth must be funded" rule above). Still the room the move (b) restructure has to spend.
%   RIGHT column : takeaway (NOW INCLUDING the 2x2 matrix) -> consequences -> model
% The matrix is no longer a separate box: it is the third movement INSIDE the Key-results card
% (block-takeaway.tex), so that card now holds highlights + mechanism cards + matrix. Its height is
% DERIVED from the logged "[layout] takeaway content height" (78.8 cm) + 22 mm of card padding.
% Bottom = 158.6 - 81.0 = 77.6, i.e. 5.7 cm above the Predictions row at \yConseq -- that clearance is
% the whole budget the card has to grow into. Re-read the log line after ANY edit to the card.
\def\yTake{158.6}\def\hTake{81.0}            % 158.6 .. 77.6 (highlights + mechanism cards + matrix)
% RETIRED with the merge: the matrix's own north edge / vertical budget.
% \def\yAbst{108.4}\def\hAbst{34}            % 2x2 matrix, when it was placed separately below the card
\def\yConseq{71.9}\def\hConseq{37}           % predictions | interpretations row
\def\yModel{50.4}                            % model-comparison card (content-sized; bottom ~20,
                                             % i.e. ~13 cm above the footer divider)
%   FULL width : footer
% (The former "Supporting details" band now lives in the standalone supplement.tex; the poster
%  points to it via the QR tag \seebottom. See HANDOFF.md.)
\def\yFoot{12}\def\hFoot{9}                   % 12 .. 3

% Coordinate helper: \xy{x}{y} -> point at (x cm, y cm) from the page corner (O).
\newcommand{\xy}[2]{($(O)+(#1 cm,#2 cm)$)}

% ===== Card (panel) styles ========================================================================
% UX-inspired cards: subtle surface contrast + soft shadow, no loud frame.
\tcbset{
  % shadow-free base: surface, rounded corners, padding.
  cardbase/.style={
    enhanced, arc=15mm, outer arc=15mm,
    colback=surface, colframe=bordersubtle, boxrule=0.8pt,
    boxsep=0pt, left=12mm, right=12mm, top=11mm, bottom=11mm,
    fontupper=\color{textprimary}\bodysize,
  },
  % homogeneous soft shadow across ALL cards: light from the front-left corner, so a single
  % tight blur is cast down and to the right (reduced spread, one direction everywhere).
  softshadow/.style={fuzzy shadow={2.5mm}{-3mm}{-0.5mm}{3.5mm}{cardshadow}},
  % deeper cast of the SAME light, for the two emphasis boxes (see emphcard).
  strongshadow/.style={fuzzy shadow={3.5mm}{-4.5mm}{-0.5mm}{5.5mm}{cardshadow}},
  card/.style={cardbase, softshadow},
  % Emphasis card (Open Question, Key results). Emphasis used to be a turquoise->white GRADIENT fill;
  % the gradient is retired. A card now earns emphasis by RELIEF, not by tint: a pronounced accent
  % ring + a deeper shadow lift it off the page while its interior stays white, so the text on it
  % reads exactly like the text on every other card.
  % TO RESTORE A GRADIENT: the two call sites keep the old fill commented out --
  % \yQuest's box in block-starting.tex (`interior style`) and the title band in block-takeaway.tex
  % (`underlay`). Uncomment there; do not put a gradient back into this style.
  emphcard/.style={cardbase, strongshadow, colback=surface, colframe=accentturq, boxrule=3pt},
  % Same emphasis card, but the shadow is TINTED BY THE RING instead of neutral grey: the box then
  % reads as glowing in its own accent rather than merely floating. Used by the Open Question, the one
  % box the reader must not miss. (Do not spread it: if every card glows, none does.)
  emphcardglow/.style={cardbase, colback=surface, colframe=accentturq, boxrule=3pt,
    fuzzy shadow={3.5mm}{-4.5mm}{-0.5mm}{6mm}{accentturq!45}},
  plaincard/.style={
    enhanced, arc=15mm, outer arc=15mm,
    colback=posterbg, colframe=posterbg, boxrule=0pt,
    boxsep=0pt, left=10mm, right=10mm, top=8mm, bottom=8mm,
    fontupper=\color{textprimary}\bodysize,
  },
  % transparent panel: same padding/alignment as a card, but no fill, frame or shadow
  % (content lies directly over the page background)
  ghostcard/.style={
    enhanced, frame hidden, interior hidden,
    boxsep=0pt, left=12mm, right=12mm, top=11mm, bottom=11mm,
    fontupper=\color{textprimary}\bodysize,
  },
  % dashed outline box (no fill, no shadow): frames the two self-consistent-gain items
  dashbox/.style={enhanced, interior hidden, frame hidden,
    arc=5mm, outer arc=5mm, boxsep=0pt, left=4mm, right=4mm, top=5mm, bottom=5mm,
    borderline={1.1pt}{0pt}{texttertiary!55, dashed}},
}

% A tcolorbox typeset directly inside a TikZ \node triggers an infinite NFSS
% font-selection loop under tikzposter + fontspec/LuaLaTeX. The fix is two-step:
% build the card into a savebox at document top level (OUTSIDE any tikzpicture),
% then place that box with \node{\usebox} inside the overlay picture.
%
% Build card CONTENT into <savebox> (call at document top level):
%   \buildcard{<\savebox>}{<width cm>}{<height cm>}{<tcb style>}{<content>}
\newcommand{\buildcard}[5]{%
  \sbox#1{\begin{tcolorbox}[#4, width=#2 cm, height=#3 cm]#5\end{tcolorbox}}%
}
% Build arbitrary (card-free) content into <savebox> at document top level:
%   \buildbox{<\savebox>}{<content>}
\newcommand{\buildbox}[2]{\sbox#1{#2}}
% Build a fixed-width text column (minipage) into <savebox>:
%   \buildcol{<\savebox>}{<width cm>}{<content>}
\newcommand{\buildcol}[3]{\sbox#1{\begin{minipage}[t]{#2 cm}#3\end{minipage}}}

% Drop a pre-built box with its NORTH-WEST corner at (x,y) cm (call INSIDE the picture):
%   \dropbox{<\savebox>}{<x>}{<y>}
\newcommand{\dropbox}[3]{%
  \node[anchor=north west, inner sep=0pt] at \xy{#2}{#3} {\usebox#1};%
}
% As \dropbox but with an explicit anchor (default north west):
%   \dropboxa[<anchor>]{<\savebox>}{<x>}{<y>}
\newcommand{\dropboxa}[4][north west]{%
  \node[anchor=#1, inner sep=0pt] at \xy{#3}{#4} {\usebox#2};%
}

% ===== Icon badge + block heading =================================================================
% Circular accent badge holding a FontAwesome glyph (or any content).
%   #1 fill color  #2 glyph color  #3 glyph
\newcommand{\iconbadge}[3]{%
  \tikz[baseline=-0.42ex]{%
    \node[circle, fill=#1, text=#2, minimum size=20mm, inner sep=0pt,
          font=\fontsize{34}{34}\selectfont]{#3};}%
}

% Functional block heading: a soft rounded-square icon (light accent tint, dark accent glyph) wearing
% its SECTION NUMBER on a small solid disc at the top-right corner, + a MONO accent kicker + a
% title-font descriptive sentence.
%   [#1 sentence font, optional]  #2 accent  #3 glyph  #4 SECTION NUMBER  #5 KICKER  #6 sentence
% The number is a reading ORDER, not a decoration: a poster is read out of order unless it says
% otherwise. It rides ON the icon's corner (a disc centred on `bh.north east') rather than sitting
% beside it, so the pair reads as one mark -- and so a card's identity is still one object, not two.
% Numbering is 1..6 in reading order: objectives, approach, key results, predictions, interpretations,
% model comparison. Renumbering means editing all six call sites; there is no central list.
\newcommand{\blockheading}[6][\fontsize{52}{60}\selectfont]{%
  \noindent
  \raisebox{-0.5\height}{%
    \tikz{%
      \node[fill=#2!16, rounded corners=9mm, minimum size=37mm, inner sep=0pt, text=#2,
        font=\fontsize{57}{57}\selectfont] (bh) {#3};
      \node[circle, fill=#2, text=white, minimum size=17mm, inner sep=0pt,
        font=\headingfont\bfseries\fontsize{30}{30}\selectfont] at (bh.north east) {#4};}}%
  \hspace{9mm}%
  \begin{minipage}[c]{\dimexpr0.9\linewidth-50mm}
    {\kickerfont\fontsize{46}{52}\selectfont\color{#2}\MakeUppercase{#5}}\par
    \vspace{1.5mm}
    {\headingfont#1\color{textprimary}#6}\par
  \end{minipage}%
  \vspace{12mm}\par
}

% Shared spacing between consecutive bullets (highlights, predictions, interpretations, ...).
\newcommand{\bulletgap}{18mm}

% Shared white space AFTER a body paragraph, before the figure row it introduces. One convention for
% every panel: a paragraph must not sit tight against the figure it announces.
% NOTE: the space after a block heading is NOT set here -- it is \blockheading's own trailing
% \vspace, and no panel may add to it (that is the convention the Predictions / Interpretations cards
% set: heading, then content, nothing in between).
\newcommand{\paragap}{12mm}

% ===== Inline citation + supporting-detail reference tag ==========================================
% Literature citation embedded at point of use, small. It takes the CHROME turquoise -- the SAME
% colour as the \seebottom supporting-detail tags below -- because the two are the same kind of mark:
% a pointer AWAY from the poster (to the literature, to the supplement). They are deliberately not
% two colours. The former semantic coral (citecoral) is retired in colors.tex.
\newcommand{\litcite}[1]{{\color{accentturq}\fontsize{32}{36}\selectfont (#1)}}
% Small numbered badge (turquoise rounded square) matching the numbered points in supplement.tex.
% The badge is placed
% with \raisebox (a LaTeX box lift) around a picture whose baseline is pinned to the node CENTRE via
% a node anchor -- neither depends on the surrounding picture's axes, so the badge aligns identically
% in ordinary text and inside a flipped-axis timeline node (e.g. the Approach, x=1cm y=-1cm).
\newcommand{\refnum}[1]{\raisebox{3.2mm}{\tikz[baseline=(rn.center)]{\node(rn)[fill=accentturq,
  text=white, rounded corners=3mm, minimum size=12mm, inner sep=0pt,
  font=\headingfont\fontsize{34}{34}\selectfont]{#1};}}}
% Pointer(s) to numbered item(s) in the standalone supplement: a small QR-code glyph + the item
% number(s), signalling "scan the header QR to reach supplemental item N". The number badges are kept.
\newcommand{\seebottom}[1]{{\color{accentturq}\fontsize{32}{36}\selectfont \mbox{\icQR\;\refnum{#1}}}}
\newcommand{\seebottomtwo}[2]{{\color{accentturq}\fontsize{32}{36}\selectfont \mbox{\icQR\;\refnum{#1}\,\refnum{#2}}}}

% ===== TikZ arrow / connector styles ==============================================================
\tikzset{
  flowarrow/.style={-{Stealth[length=19pt,width=15pt]}, line width=2.2pt,
    rounded corners=8pt},
  softarrow/.style={-{Stealth[length=15pt,width=11pt]}, line width=1.6pt},
  thinarrow/.style={-{Stealth[length=11pt,width=8pt]}, line width=1.0pt},
  guideflow/.style={-{Stealth[length=22pt,width=16pt]}, line width=3pt,
    accentturq!55, rounded corners=14pt},
  % A COORDINATE MARK: the point on a vector at which a coefficient is read off (alpha_k u_k, beta_k k_k,
  % kappa_k m_k). It sits ON the vector's shaft, so a plain \fill of the same colour disappears into it --
  % the WHITE RING is what makes it a mark rather than a thicker piece of the arrow. The ring's line width
  % is absolute (TikZ does not scale line widths), so the mark keeps the same rim in every figure; the
  % RADIUS, however, IS scaled by the picture's `scale', so each figure passes a radius that lands at
  % ~4.4pt on the page (2.4pt at scale 1.85, 4pt at scale 1.1, ...). Use as: \draw[coordmark, fill=compred].
  coordmark/.style={draw=white, line width=1.2pt},
}
